% ============================================================
% REVISED SECTION 7: DISCUSSION
% Reorganized into three subsections (7.1, 7.2, 7.3) and
% expanded to address (i) entry-rate dynamics and (ii) the
% evolution of z* under the wage channel.
% ============================================================

    The findings in the previous section point to increases in wages, firm size, entry, and exit rates following an exogenous immigration shock. These findings are consistent with increased business dynamism and concentration caused by the arrival of immigrants. I use the model developed in Section \ref{sec:model} to explain the forces behind these changes and to infer the type of these arriving immigrants. However, the model is a partial equilibrium, single-period framework in which the wage $w$ is taken as exogenous and flow concepts such as entry and exit rates are not explicitly derived. Therefore, the discussion below distinguishes carefully between results that follow from the model's equations and interpretations that follow from its economic intuition. The limitations subsection at the end consolidates these caveats.

\subsection{The Wage Channel and Managerial Talent}


    I will use economic intuition to argue that wages can be tied directly to managerial talent. A simple extension of the model can pin down the equilibrium wage given stationary distributions of incumbents and managerial talents. For now, I focus on Equation~\eqref{eq:wage}, which shows the relation between wages and the managerial talent cut-off for new entrants. One can see that there is a positive relationship between wages and managerial talent: an increase in $w$ shifts the unconstrained cutoff $\theta^{*}_{u}$ proportionally upward, raising the level of managerial talent one must have for entrepreneurship to dominate the worker outside option. If one observes an increase in wages, this must be accompanied by an increase in the prevailing talent cut-off. These effects are consistent with the story of arriving immigrants being more talented than the existing population.

    I can now paint a complete picture of the effects in Table~\ref{tab:main}. Column (1) shows an increase in wages, which is interrelated with increased managerial talent in the economy. Column (2) shows increased firm size, a direct consequence of more talented managers demanding more labor on
    average (see Equation~\eqref{eq:labordemand}). Columns (3) and (4) show increases in both the entry rate and the exit rate, respectively. While the model makes no formal predictions about transition dynamics, the overall findings are consistent with the arrival of higher-talent migrants who set up new firms while older, less-talented incumbents exit. I unpack the entry and exit margins separately in the next subsection.

\subsection{Entry, Exit, and the Model's Comparative Statics}

    \paragraph{Entry rate:} The positive entry-rate response in
    Column (3) of Table \ref{tab:main} deserves further attention, as the model's entry condition $V^{E}(n,k;w)\geq w$ is informative about its sign even though the model itself is single-period and cannot speak to flow rates directly. Two channels operate simultaneously when positively selected immigrants arrive. First, arriving immigrants with high $V^{E}_{I}$ enter directly at rates above the prevailing native average, mechanically raising the county-level entry rate. Second, the induced increase in $w$ tightens the entry condition for natives near the margin: some types $(N,k)$ that previously satisfied $V^{E}_{N}(N,k;w_{\text{old}}) \geq w_{\text{old}}$ no longer satisfy $V^{E}_{N} (N,k;w_{\text{new}}) \geq w_{\text{new}}$. The empirical finding of a positive net entry rate response indicates that the direct-entry channel dominates this crowding-out of marginal native entrants. The model does not pin down the relative magnitude of these two channels, but its comparative-statics intuition rationalizes both the sign and the broader pattern of more churn at higher equilibrium wages.

    \paragraph{Exit rate and the evolution of $z^{*}$:} Table~\ref{tab:bdsage} provides a clearer view of the exit margin by stratifying by firm age. Panel~A shows that the increased exit rate is driven by medium and old firms (those 6+ years old), not by new arrivals. Since the empirical analysis is conducted in 5-year periods, firms in columns~(3) and~(4) of Table~\ref{tab:bdsage} were established before the immigration shock arrived. The increase in exit rate is therefore not driven by failed immigrant start-ups but reflects a broader general-equilibrium effect on incumbent businesses.

    The model's exit cutoff $z^{*}(\theta,k,w)$ in Equation~\eqref{eq:z-star} provides intuition for this pattern. Recall that $z^{*}$ is strictly increasing in $w$ and strictly decreasing in $\theta$: higher wages raise the threshold below which a transitory shock forces an incumbent to exit, and this threshold disproportionately affects the low-$\theta$ incumbents. When immigration raises $w$, every incumbent's $z^{*}$ rises, and the per-period exit probability $d(\theta,k,w)=H(z^{*}(\theta,k,w))$ rises most for those operating at low realized talent. The exiting incumbent firms are precisely those who have survived past entry but were operating at margins of the $z^*$ cut-off. When $w$ rises, these marginal incumbents are forced out. The model is single-period and cannot formally derive the cross-period evolution of $z^{*}$, but its comparative statics rationalize the empirical finding that exit rates rise mainly for medium and old firms following an immigration shock.

    \paragraph{Firm size across ages:} Panel~B of Table~\ref{tab:bdsage} shows that the increase in firm size is spread roughly equally across firm age buckets. The estimates for the four age groups are similar to the overall coefficient in Table~\ref{tab:main}, with no apparent heterogeneity. This is consistent with a story in which the average managerial talent of operating firms has increased across all ages: among new firms, through the direct entry of higher $\theta$ immigrant entrepreneurs; among older firms, through the selection mechanism described above, wherein the rising $z^{*}$ removes low $\theta$ incumbents and leaves a more positively selected pool. In both cases, the surviving firms hire more labor on average, consistent with Equation \eqref{eq:labordemand}.

    \paragraph{Summary:} My empirical findings are consistent with the idea that arriving migrants are more talented than the existing population in a county. These immigrants raise the equilibrium wage and through it, the talent cutoffs $\theta^{*}$ for entry and $z^{*}$ for incumbent survival. Marginal incumbents exit, new high talent immigrants enter, and firm size rises across the board. All of these effects are mediated via the wage channel, which the model treats as exogenous but whose role is transparent in the entry and exit conditions.

\subsection{Limitations and Scope of Interpretation}


    The findings of this study should be interpreted keeping in mind several limitations and simplifying assumptions. Several of these arise because the empirical findings are interpreted through the model's economic intuition rather than through formally derived equations. 
    
    First, the model is single period, so flow concepts like entry and exit rates are not directly defined within it. The discussion above invokes the entry condition $V^{E}\geq w$ and the exit cutoff $z^{*}$ as comparative-static objects rather than as dynamic equilibrium quantities. The cross-period evolution of $z^{*}$ following an immigration shock, including the transition path of the incumbent distribution and the eventual stationary distribution, is outside the model's scope.
    
    Second, the wage response itself is exogenous to the model. The channels through which immigration shifts $w$, and subsequently $V^{E}$ and $z^{*}$, are described informally. A natural extension would embed the framework in a dynamic general-equilibrium setting with an endogenous wage and a stationary distribution of incumbents, allowing the entry-rate, exit-rate, and $z^{*}$ comparative statics to be derived formally rather than invoked as intuition. I leave this for future work.

    Third, I do not account for any productivity growth associated with migration. Recent work has highlighted the importance of these effects and shown that immigration is strongly connected to productivity growth in the US (\citealt{EffectofImmigrationProductivity}; \citealt{ImmigrationInnovationGrowth}). My partial-equilibrium framework holds the production technology fixed and attributes all cross-county variation in firm-size and wage outcomes to the talent channel. In other words, the gains for a local economy induced by immigration may be larger than the ones shown in this study.

    Fourth, I do not account for any correlation or signal between capital endowment and managerial talent. It is highly plausible that these are positively correlated, and incorporating this would alter the predictions about which $(n,k)$ types enter entrepreneurship. Future work can try to overcome some of these limitations to provide a more robust understanding of immigrant entrepreneurship in the United States.

% ============================================================
% End of revised Section 7.
% ============================================================